\documentclass{article}
\usepackage[T1]{fontenc}
\usepackage{lmodern}
\usepackage{graphicx}
\usepackage{amsmath,amssymb}
\usepackage[a4paper,margin=2cm]{geometry}
\usepackage[absolute,overlay]{textpos}
\setlength{\TPHorizModule}{1mm}
\setlength{\TPVertModule}{1mm}
\usepackage[indonesian]{babel}
\usepackage{hyperref}
\setlength{\emergencystretch}{2em}
% unit: Halaman 43

\begin{document}

% === Halaman 43 (python/native) ===

Contoh *):  Misalkan kurva eliptik yang dipilih adalah y2 = x3 + 2x + 1 mod  5. Himpunan 

titik-titik pada kurva eliptik adalah \{(0, 1), (1, 3), (3, 3), (3, 2), (1, 2), (0, 4)\}. Alice dan Bob 

menyepakati titik B(0, 1) sebagai basis.

1.

Alice memilih a = 2, lalu menghitung kunci publiknya:



PA = aB = 2B = B + B = (1, 3)  → misalkan titik Q

2.

Bob memilih b = 3, lalu menghitung kunci publiknya:



PB = bB = 3B = B + B + B = 2B + B = (3, 3) → misalkan titik R

3.

Alice mengirimkan PA kepada Bob, Bob mengirimkan PB kepada Alice.

4.

Alice menghitung kunci rahasia sbb:



KA = aPB = 2R = R + R = (0, 4)

5.

 Bob menghitung kunci rahasia sbb:



KB = bPA = 3Q = Q + Q + Q = (0, 4)

Jadi, sekarang Alice dan Bob sudah berbagi kunci rahasia yang sama, yaitu (0, 4)


43

*) Sumber bahan: Nana Juhana, Implementasi Elliptic Curve 

Cryptography  (ECC) pada proses Pertukaran Kunci Diffie-Hellman 

dan Skema Enkripsi El Gamal

\end{document}
